% BHU PYQ -- Manually transcribed & error-corrected
% Paper: GLB-503: Sedimentary and Metamorphic Petrology
% Exam: B.Sc. (Hons.) Semester V Examination, 2016-17

\documentclass[11pt,a4paper]{article}
\usepackage[a4paper,top=2.5cm,bottom=2.5cm,left=2.8cm,right=2.8cm,headheight=14pt]{geometry}
\usepackage{amsmath,amssymb}
\usepackage[T1]{fontenc}
\usepackage[utf8]{inputenc}
\usepackage{lmodern,microtype,fancyhdr,enumitem,hyperref,lastpage,parskip}

\hypersetup{colorlinks=true,urlcolor=black,linkcolor=black,
  pdftitle={GLB-503 Sedimentary and Metamorphic Petrology -- BHU B.Sc. Sem V 2016-17}}
\pagestyle{fancy}\fancyhf{}
\renewcommand{\headrulewidth}{0.4pt}\renewcommand{\footrulewidth}{0.4pt}
\fancyhead[L]{\small\textit{Banaras Hindu University}}
\fancyhead[R]{\small\textit{GLB-503: Sedimentary and Metamorphic Petrology}}
\fancyfoot[C]{\small Page \thepage\ of \pageref{LastPage}}

\newlist{parts}{enumerate}{2}
\setlist[parts,1]{label=(\alph*),leftmargin=*,itemsep=5pt,topsep=3pt}
\newcommand{\pts}[1]{\hfill\textit{[#1]}}

\begin{document}
\begin{center}
  {\large\bfseries BANARAS HINDU UNIVERSITY}\\[2pt]
  {\normalsize B.Sc. (Hons.) --- Semester V Examination, 2016-17}\\[6pt]
  \rule{0.8\linewidth}{0.5pt}\\[6pt]
  {\large\bfseries Paper No. GLB-503: Sedimentary and Metamorphic Petrology}\\[4pt]
  {\normalsize Time: 3 Hours \hfill Maximum Marks: 70}\\[2pt]
  {\small Ref. No.: BS/Sem V/174}
\end{center}
\rule{\linewidth}{0.4pt}
\smallskip
\noindent\textit{Instructions: Attempt any \textbf{five} questions, selecting at least \textbf{two} questions each from Section A and B. Section C is compulsory. All questions carry equal marks (14 marks each).}
\bigskip

\bigskip\noindent{\large\bfseries SECTION --- A}
\rule{\linewidth}{0.4pt}\smallskip

\noindent\textbf{Question 1.} Discuss in detail the processes responsible for the formation of sedimentary rocks. \pts{14}

\medskip\noindent\textbf{Question 2.} Describe the classification of conglomerates and characterize their genesis. \pts{14}

\medskip\noindent\textbf{Question 3.} Write detailed notes on \textbf{any two} of the following: \pts{$2 \times 7 = 14$}
\begin{parts}
  \item Argillaceous rocks
  \item Cross-bedding
  \item Nonclastic texture
  \item Heavy minerals
\end{parts}

\bigskip\noindent{\large\bfseries SECTION --- B}
\rule{\linewidth}{0.4pt}\smallskip

\noindent\textbf{Question 4.} Describe progressive metamorphism of ultramafic rocks in the $\text{MgO}-\text{Al}_2\text{O}_3-\text{SiO}_2-\text{H}_2\text{O}$ system. \pts{14}

\medskip\noindent\textbf{Question 5.} Discuss the metamorphic facies and metamorphic facies series. \pts{14}

\medskip\noindent\textbf{Question 6.} Write notes on \textbf{any two} of the following: \pts{$2 \times 7 = 14$}
\begin{parts}
  \item Goldschmidt's mineralogical Phase Rule
  \item Prograde Metamorphism
  \item ACF and AKF diagrams
\end{parts}

\medskip\noindent\textbf{Question 7.} Describe progressive metamorphism of Pelitic rocks in the $\text{K}_2\text{O}-\text{FeO}-\text{MgO}-\text{Al}_2\text{O}_3-\text{SiO}_2$ system. \pts{14}

\bigskip\noindent{\large\bfseries SECTION --- C (Compulsory)}
\rule{\linewidth}{0.4pt}\smallskip

\noindent\textbf{Question 8.} Fill in the blanks: \pts{$7 \times 2 = 14$}
\begin{parts}
  \item Eclogites are composed of \underline{\hspace{5cm}} and \underline{\hspace{5cm}}.
  \item Quartzarenites are texturally \underline{\hspace{3cm}} and mineralogically \underline{\hspace{3cm}}.
  \item In the series Slate --- Phyllite --- Schist, the grain size of the rock \underline{\hspace{4cm}}.
  \item Small-scale ripples have a height of \underline{\hspace{4cm}}.
  \item $\text{Hornblende} + \text{quartz} = \text{\underline{\hspace{4cm}}} + \text{diopside} + \text{plagioclase} + \text{H}_2\text{O}$.
  \item Freshly deposited clay has a porosity of \underline{\hspace{4cm}}.
  \item $\text{Corundum} + \text{quartz} = \text{\underline{\hspace{4cm}}}$ (Choose from: \textit{sillimanite / sapphirine / cordierite}).
\end{parts}

\vfill\rule{\linewidth}{0.4pt}
\begin{center}\small
\href{https://bhu-exam.vercel.app/}{Click here to visit our website}\\[4pt]
\href{https://me-aryan.vercel.app/}{Click here to visit our contributor}\\[4pt]
\href{https://quashan-me.vercel.app/}{Click here to visit our contributor}
\end{center}
\end{document}
